\section{Version history}\label{app:history}

\paragraph{v2 (October 1, 2026; DOI \href{https://doi.org/10.5281/zenodo.23088986}{10.5281/zenodo.23088986}).} The currency--constraint principle is stated as a conditional theorem, with the resource law as an explicit premise and a counterexample showing that persistence alone does not supply one. The finite Gibbs duality is proved in full, with all price regimes, and its single-row form is mechanized in Lean. The data-processing theorem is extended to arbitrary deterministic maps and sparse supports, with a path-reversal form. Empirical path audits treat unmatched reversals as infinite asymmetry and use a reversal-mixture regularization that satisfies data processing on the same sample. The proxy experiment uses a fixed design and a fixed proxy. The manuscript uses held-out prediction as the supported comparison and no longer draws a price-stability conclusion: across-sample price stability, reported as a comparison in v1, is not supported under the fixed design. The experiment still records the price dispersions and the outcome of its original criterion. The ladder result separates available cycle dimension from driven dimension. A uniform packaging-defect bound, and the distinction between calibration cost and actual spending, were added. Figures and the exposition were rewritten throughout.

\paragraph{v1 (March 9, 2026; DOI \href{https://doi.org/10.5281/zenodo.18926771}{10.5281/zenodo.18926771}).} First public version.
